\documentclass[11pt]{amsart}
% packages
\usepackage[margin=1.1in]{geometry}
\usepackage{amssymb,amsthm}
\usepackage{xcolor}
\usepackage[colorlinks=true,linkcolor=blue!50!black,citecolor=blue!50!black,urlcolor=blue!50!black]{hyperref}

% theorem styles
\theoremstyle{plain}
	\newtheorem{theorem}{Theorem}
	\newtheorem{lemma}[theorem]{Lemma}
	\newtheorem{corollary}[theorem]{Corollary}
\theoremstyle{remark}
	\newtheorem{remark}[theorem]{Remark}

% shortcuts
\newcommand{\Mlt}{\mathrm{Mlt}}
\newcommand{\Stab}{\mathrm{Stab}}

\newcommand{\E}{\mathbb{E}}
\renewcommand{\Pr}{\mathbb{P}}
\newcommand{\LL}{\mathcal{L}}
\DeclareMathOperator{\sgn}{sgn}

\title{Almost all loops have multiplication group $S_n$}
\author{}
\date{October 2026 --- draft}

\begin{document}
\maketitle

\begin{abstract}
Cameron proved that for almost all quasigroups of order $n$ the left
translations generate the symmetric or the alternating group, and
H\"aggkvist and Janssen excluded the alternating group; Cameron asked
whether the same holds for loops, where the argument breaks because the
translation by the identity is trivial.  We show that it does, by counting: the
probability that the rows of a uniformly random reduced Latin square of
order $n$ generate a group not containing $A_n$ is at most
$e^{-(\log 2-o(1))n^2}$, and the group is $S_n$ with probability
$1-O(c^n)$ for some $c<1$.  The inputs are the van der Waerden and
Br\'egman bounds on the number of Latin squares, the Praeger--Saxl bound
on the orders of primitive groups, Pyber's bound on the number of
subgroups of $S_n$, and the H\"aggkvist--Janssen bound on Latin squares
all of whose rows are even permutations.  This settles the question,
raised by Cameron and recorded by Cavenagh, Greenhill and Wanless as a
consequence of their conjecture on the second row, and completes an
argument sketched by Eberhard in 2016.
\end{abstract}

\section{Introduction}

Let $Q$ be a quasigroup of order $n$ on the set $[n]=\{1,\dots,n\}$,
with Cayley table a Latin square $L$, and for $x\in Q$ let
$L_x\colon y\mapsto xy$ be the left translation, the $x$-th row of $L$
read as a permutation of $[n]$.  The \emph{left multiplication group}
$\Mlt_\lambda(Q)=\langle L_x:x\in Q\rangle$ is a transitive subgroup
of $S_n$, and the full multiplication group $\Mlt(Q)$, generated by all
left and right translations, contains it.  Cameron \cite{Cam92} showed
that for almost all quasigroups (that is, for all but a vanishing
proportion of the Latin squares of order $n$) the left translations
generate $S_n$ or $A_n$; H\"aggkvist and Janssen \cite{HJ96} then showed
that the proportion of Latin squares all of whose rows are even
permutations is at most $c^n$ for some $c<1$, so that $A_n$ never
occurs.  In particular $\Mlt(Q)$ is $2$-transitive, which in Smith's
character theory of quasigroups \cite{Smith} means that almost all
quasigroups have trivial character theory (``rank $2$'').

Cameron's argument is short.  A uniformly random Latin square has a
uniformly random first row, and by a theorem of {\L}uczak and Pyber
\cite{LP93} a uniformly random permutation lies in no transitive proper
subgroup of $S_n$ other than $A_n$ with probability $1-o(1)$; since the
rows generate a transitive group, that group contains $A_n$.

For \emph{loops} --- quasigroups with an identity element, whose Cayley
tables, suitably labelled, are the \emph{reduced} Latin squares, with
first row and first column in natural order --- the argument is not
available: the first row is the identity permutation, and no other row
is uniformly distributed.  Cameron asked whether almost all loops nevertheless have
multiplication group $S_n$ (see \cite{Cam15}, and \cite[\S6]{CGW08},
where Cavenagh, Greenhill and Wanless record, as Lemma~6.2, that this
follows from their conjecture on the cycle structure of the second row
of a random Latin square).  The statement is folklore in loop theory,
where it is mentioned as an unproved belief.  The purpose of this note
is to settle it, by an argument that uses no information about the
second row at all.

\begin{theorem}\label{thm:main}
Let $L$ be a uniformly random reduced Latin square of order $n$ and
$H(L)\le S_n$ the group generated by its rows.  Then
\[
\Pr\bigl[H(L)\not\supseteq A_n\bigr]\le e^{-(\log2-o(1))n^2},
\qquad\text{and}\qquad
\Pr\bigl[H(L)=S_n\bigr]\ge1-3c^n
\]
for some absolute $c<1$ and all large $n$.  Consequently almost all loops
have $\Mlt_\lambda(Q)=\Mlt(Q)=S_n$ and trivial character theory.
\end{theorem}

The idea, which Eberhard sketched in comments on Cameron's post
\cite{Cam15}, is that the event is about all $n-1$ non-identity rows
simultaneously, and such events are cheap to bound by counting: for a
subgroup $K\le S_n$, at most $|K|^{n-1}$ of the $L(n)/n!=(n!)^{n-1}e^{-n^2+o(n^2)}$
Latin squares with first row the identity have all their rows in $K$,
a proportion of at most $e^{n^2+o(n^2)}/[S_n:K]^{n-1}$, which is
$e^{-\delta n^2+o(n^2)}$ as soon as $[S_n:K]\ge e^{(1+\delta)n}$.  A
transitive proper subgroup lies in a maximal one, and the maximal
transitive subgroups of $S_n$ are the wreath products $S_a\wr S_b$,
$ab=n$, of which there are fewer than $n\cdot n!$, and the primitive
maximal subgroups.  The crude count disposes of $S_a\wr S_b$ for
$b\ge3$ (index at least $3^n/n^{O(1)}$) but not of $S_{n/2}\wr S_2$
(index $\binom n{n/2}/2<e^n$) nor of $A_n$; for the primitive
subgroups one needs a bound on their orders and on their number.  We
handle every imprimitive group by a single block decomposition
(Lemma~\ref{lem:block}), which is sharper than the crude count ---
it gives the exponent $\log2$ in place of $\log3-1$ --- and in
particular disposes of $S_{n/2}\wr S_2$; the primitive groups by the
Praeger--Saxl order bound together with Pyber's bound on the number of
subgroups of $S_n$; and $A_n$ by the theorem of H\"aggkvist and
Janssen.

\section{Preliminaries}

\paragraph{Latin squares.}  Let $\LL_n$ be the set of Latin squares of
order $n$ on the symbol set $[n]$ and $L(n)=|\LL_n|$.  We use two
classical bounds \cite[Ch.~17]{vLW}: the lower bound
$L(n)\ge(n!)^{2n}/n^{n^2}$, a consequence of the van der Waerden
permanent inequality applied row by row, and the upper bound
$L(n)\le\prod_{k=1}^n(k!)^{n/k}$, a consequence of Br\'egman's.  In the
explicit forms we need,
\begin{equation}\label{eq:L}
\log L(n)\ge n^2\log n-2n^2,\qquad
\log L(a)\le a^2\log a-2a^2+a\bigl(\tfrac12\log^2a+2\log a+3\bigr)
\end{equation}
for all $n,a\ge1$.  The first follows from $n!\ge(n/e)^n$.  For the
second, $\log k!\le k\log k-k+\log k+1$ for $k\ge1$ gives
$(\log k!)/k\le\log k-1+(\log k+1)/k$, and summing over $k\le a$ with
$\log a!\le a\log a-a+\log a+1$, $\sum_{k\le a}(\log k)/k\le\frac12\log^2a+1$
and $\sum_{k\le a}1/k\le\log a+1$,
\[
\sum_{k=1}^a\frac{\log k!}k\le a\log a-2a+\tfrac12\log^2a+2\log a+3,
\]
which multiplied by $a$ is the claim.

\paragraph{Loops and reduced squares.}  For $L\in\LL_n$ with rows
$r_1,\dots,r_n\in S_n$ (so $r_i(c)=L(i,c)$), relabelling the symbols by
$r_1^{-1}$ gives the square with rows $\sigma_i=r_1^{-1}r_i$, whose
first row is the identity; this map $\LL_n\to\LL_n^0:=\{L:r_1=\mathrm{id}\}$
is $n!$-to-one, so a proportion of squares in $\LL_n^0$ is the same proportion
of squares in $\LL_n$, and $|\LL_n^0|=L(n)/n!$.  Permuting the rows of $L^0$ so that its first
column is $1,2,\dots,n$ gives a reduced square; each reduced square
arises from exactly $(n-1)!$ squares $L^0$, so proportions agree again,
and the two have the same \emph{set} of rows.
A uniform reduced square is the Cayley table of a uniform random loop
on $[n]$ with identity $1$, whose left translations are its rows.
Writing
\[
H(L):=\langle\sigma_2,\dots,\sigma_n\rangle=\langle r_1^{-1}r_i:2\le i\le n\rangle\le S_n,
\]
it follows that for every subgroup $K\le S_n$, the proportion of $L\in\LL_n$ with $H(L)=K$ equals the proportion of loops $Q$ of order $n$ with $\mathrm{Mlt}_\lambda(Q)=K$, and it suffices to prove Theorem~\ref{thm:main} for $H(L)$ with $L$ uniform on $\LL_n$ (or, equivalently, $L^0$ uniform on $\LL_n^0$).
$H(L)$ is transitive: for $c,s\in[n]$ the symbol $r_1(s)$ occurs in column $c$ in some row $i$,
and then $\sigma_i(c)=s$.

\paragraph{Group theory.}  We use three facts about subgroups of $S_n$.
A transitive subgroup not containing $A_n$ is either imprimitive, i.e.\
contained in the stabiliser $\Stab(\Pi)\cong S_a\wr S_b$ of a partition
$\Pi$ of $[n]$ into $b$ blocks of size $a$ with $1<a<n$, or primitive;
a primitive subgroup not containing $A_n$ has order less than $4^n$
(Praeger and Saxl \cite{PS80}); and the number of subgroups of $S_n$ is
at most $2^{Cn^2}$ for an absolute constant $C$ (Pyber \cite{Pyb93},
with $C=\frac16\log_224+o(1)$; Roney-Dougal and Tracey \cite{RT25} have
recently shown $2^{n^2/16+O(n^{3/2})}$).

\section{Proof of Theorem~\ref{thm:main}}

\begin{lemma}[squares with a block structure]\label{lem:block}
Let $\Pi=\{C_1,\dots,C_b\}$ be a partition of the columns $[n]$ and
$\Pi'=\{S_1,\dots,S_b\}$ a partition of the symbols $[n]$, both into $b$
blocks of size $a$, $ab=n$.  The number of Latin squares $L\in\LL_n$
every row of which maps each block of $\Pi$ onto a block of $\Pi'$ is at
most $(b!)^n\,L(a)^{b^2}$.
\end{lemma}

\begin{proof}
Each row $r_i$ of such a square induces a bijection $\pi_i\colon[b]\to[b]$
with $r_i(C_j)=S_{\pi_i(j)}$.  Fix $j,j'$ and let
$I_{jj'}=\{i:\pi_i(j)=j'\}$.  In a column $c\in C_j$ each of the $a$
symbols of $S_{j'}$ occurs exactly once, and only in rows of $I_{jj'}$,
each of which contributes exactly one symbol of $S_{j'}$ to column $c$;
hence $|I_{jj'}|=a$, and the subarray with rows $I_{jj'}$ and columns
$C_j$ is an $a\times a$ Latin square on the symbol set $S_{j'}$.  Thus
$L$ is determined by the $n\times b$ array $(\pi_i(j))$, whose rows are
permutations of $[b]$, together with one Latin square of order $a$ for
each of the $b^2$ pairs $(j,j')$, on the row set $I_{jj'}$, column set
$C_j$ and symbol set $S_{j'}$ determined by the array (identifying
each of these sets with $[a]$ in increasing order).  There are at most
$(b!)^n$ arrays and $L(a)$ choices for each small square.
\end{proof}

\begin{proof}[Proof of Theorem~\ref{thm:main}]
Throughout, $\Pr$ denotes proportion in $\LL_n^0$; by the discussion above,
these proportions equal the corresponding proportions of reduced squares, hence of loops.
Let $L^0$ be uniform on $\LL_n^0$ and $H=H(L^0)$. Since $H$ is
transitive, if $H\not\supseteq A_n$ then $H$ is imprimitive or
primitive of order less than $4^n$.

\emph{Imprimitive.}  Suppose $H$ preserves a partition $\Pi$ of $[n]$
into $b$ blocks of size $a$, $1<a<n$.  Then every row of $L^0$ maps
each block of $\Pi$ onto a block of $\Pi$, so by Lemma~\ref{lem:block}
with $\Pi'=\Pi$,
\[
\Pr\bigl[H\le\Stab(\Pi)\bigr]\le\frac{(b!)^n\,L(a)^{b^2}}{L(n)/n!}.
\]
There are $n!/((a!)^bb!)\le n!$ partitions $\Pi$ with blocks of size $a$
and fewer than $n$ divisors $a$ of $n$.  By \eqref{eq:L}, with $b=n/a$,
$(b!)^n\le b^{nb}$ and $\log n!\le n\log n$,
\begin{align*}
\log\frac{n!\cdot n!\,(b!)^n\,L(a)^{b^2}}{L(n)}
&\le2n\log n+\frac{n^2}a\log b+\Bigl(n^2\log a-2n^2+\frac{n^2}a\bigl(\tfrac12\log^2a+2\log a+3\bigr)\Bigr)-n^2\log n+2n^2\\
&=-n^2\log b+\frac{n^2}a\Bigl(\log b+\tfrac12\log^2a+2\log a+3\Bigr)+2n\log n.
\end{align*}
If $a\ge\sqrt n$, the middle term is at most
$n^{3/2}(\frac12\log^2n+3\log n+3)=o(n^2)$, so the right side is
$-n^2\log b+o(n^2)\le-n^2\log2+o(n^2)$.  If $2\le a<\sqrt n$, then
$b>\sqrt n$, $(n^2/a)\log b\le\frac12n^2\log b$, and
$(\frac12\log^2a+2\log a+3)/a\le\frac52$ (the function is decreasing in
$a\ge1$ and is less than $2.32$ at $a=2$), so the right side is at most
$-\frac12n^2\log b+\frac52n^2+2n\log n\le-\frac14n^2\log n+3n^2$, which
is at most $-n^2\log2$ once $n\ge e^{15}$.  Summing over $\Pi$ and $a$,
\[
\Pr[H\text{ imprimitive}]\le n\,e^{-(\log2-o(1))n^2},
\]
where the $o(1)$ is $O(\log^2n/n)$, from $a=n/2$.

\emph{Primitive.}  For a fixed subgroup $K\le S_n$, the squares
$L^0\in\LL_n^0$ all of whose rows lie in $K$ number at most $|K|^{n-1}$.
Summing over the at most $2^{Cn^2}$ subgroups $K$ of order less than
$4^n$,
\[
\Pr\bigl[H\text{ primitive},\ H\not\supseteq A_n\bigr]
\le\frac{2^{Cn^2}\,4^{n(n-1)}}{L(n)/n!}
\le\exp\bigl(-n^2\log n+(C\log2+\log4+2)n^2+n\log n\bigr)
=e^{-(1-o(1))n^2\log n}.
\]
Adding the two cases gives the first inequality of the theorem.

\emph{The sign.}  $H\le A_n$ if and only if every $\sigma_i=r_1^{-1}r_i$
is even, i.e.\ every row of $L$ has the parity of $r_1$.  Relabelling
the symbols by a transposition is a bijection of $\LL_n$ reversing the
parity of every row, so this event has probability
$2\Pr[\text{all rows of }L\text{ are even}]\le2c^n$ by
H\"aggkvist and Janssen \cite{HJ96}.  If $H\supseteq A_n$ and
$H\not\le A_n$ then $H=S_n$; so $\Pr[H\ne S_n]\le e^{-(\log2-o(1))n^2}+2c^n\le3c^n$
for large $n$.  Finally $\Mlt(Q)\supseteq\Mlt_\lambda(Q)=S_n$ is
$2$-transitive, which is the rank-$2$ condition of \cite{Smith}.
\end{proof}

\section{Remarks}

\begin{remark}[sharpness]
The dominant term is $S_{n/2}\wr S_2$, and there the exponent is right:
the squares all of whose rows preserve a fixed partition into two
blocks of size $n/2$ number $\binom n{n/2}L(n/2)^4$ by the bijective
version of Lemma~\ref{lem:block}, and
$\binom n{n/2}L(n/2)^4/L(n)=2^{-n^2+O(n\log^2n)}$ by the asymptotics
$\log L(m)=m^2\log m-2m^2+O(m\log^2m)$ that \eqref{eq:L} encodes
(this counts squares with arbitrary first row; the proportion among those with first row the identity is larger by $\tfrac12\binom{n}{n/2}=2^{O(n)}$, which does not affect the exponent); the
union over the $\frac12\binom n{n/2}$ such partitions changes this by
$2^{O(n)}$.  The threshold $n\ge e^{15}$ in the proof is an artefact of
the case split and has not been optimised: the displayed bound with
\eqref{eq:L} inserted is already less than $1$ for $n\ge250$ (and equals
$e^{-0.88\,n^2\log2}$ at $n=1000$), and with the exact van der Waerden
and Br\'egman expressions in place of \eqref{eq:L} it is less than $1$
for $n\ge30$ ($e^{-0.37\,n^2\log2}$ at $n=30$, $e^{-0.71\,n^2\log 2}$
at $n=100$).
\end{remark}

\begin{remark}[quasigroups]
With $G=\langle r_1,\dots,r_n\rangle$ in place of $H$, the same three
steps (Lemma~\ref{lem:block} with $\Pi'=\Pi$, the count $|K|^n$, and
\cite{HJ96}) reprove Cameron's theorem \cite{Cam92} without
\cite{LP93}, with the same error bound.  The converse direction is not
available: the loop case cannot be deduced from the quasigroup case,
because the symbol relabelling that produces a reduced square replaces
the row group $\langle r_i\rangle$ by $\langle r_1^{-1}r_i\rangle$.
\end{remark}

\begin{remark}[what is and is not used]
Nothing is used about the distribution of any single row.  The event
$H\le K$ is the intersection of the $n-1$ events $\sigma_i\in K$, and the
identity row, which defeats the one-row argument, contributes nothing
to it.  The crude count $|K|^{n-1}$ against $(n!)^{n-1}e^{-n^2+o(n^2)}$
already settles each individual $K$ of index at least
$e^{(1+\delta)n}$; the two maximal subgroups of smaller index, $A_n$
and $S_{n/2}\wr S_2$ (with their conjugates), are where the Latin
constraints inside $K$ have to be used, through \cite{HJ96} and
through Lemma~\ref{lem:block} respectively.  In particular the theorem
says nothing about the conjecture of \cite{CGW08} on the cycle type of
the second row, which concerns one row given the identity row, and
which remains open; the consequence recorded in \cite[Lemma~6.2]{CGW08}
simply does not need it.
\end{remark}

\begin{remark}[Eberhard's sketch]
In comments on \cite{Cam15} (October 2016), Eberhard observed that
the number of Latin squares with all rows in $K$ is at most
$|K|^n=(n!)^n/[S_n:K]^n$ against $L(n)=(n!)^n/e^{n^2+o(n^2)}$, concluded
that only $A_n$ and $S_{n/2}\wr S_2$ need separate treatment, asked
whether the method of \cite{HJ96} adapts to other subgroups, and
sketched for $S_{n/2}\wr S_2$ a row-by-row count (fill the first
$k\approx n/10$ rows inside the group, bound the rest by the standard
row-by-row estimate); Cameron replied that this ``should not be too
hard'' and that the second-row conjecture was the interesting part.
The sketch is correct once one passes to maximal subgroups, which
takes care of the imprimitive groups with $b\ge3$; what it leaves open
is the primitive case, which needs the order bound of \cite{PS80} and
a count of subgroups such as \cite{Pyb93}, and the details for
$S_{n/2}\wr S_2$.  Lemma~\ref{lem:block} treats all imprimitive groups
uniformly and gives the sharp exponent.  As far as we know the argument
was not written out.
\end{remark}

\begin{remark}[small orders]
For $n=5$, of the $56$ reduced squares $50$ have $H=S_5$ and $6$ have
$H=C_5$ (the sorted Cayley table of $\mathbb Z_5$ for each of the six
cyclic subgroups; these are also the all-even ones).  For $n=6$, of
$9408$ reduced squares $7776$ have $H=S_6$, $72$ have $A_6$ and $1560$
have an imprimitive $H$ (block sizes $2$ and $3$, orders $6$ to $48$).
In samples of uniformly random Latin squares of orders $7$ to $16$
(Jacobson--Matthews chain, $2000$--$3000$ squares per order) every $H$
was $S_n$ or $A_n$, with $A_n$ in $2.1\%$, $1.2\%$, $0.4\%$, $0.1\%$ of
the samples at $n=7,8,9,10$ and never at $n=12,16$.
\end{remark}

\begin{remark}[labelled and unlabelled loops]
The theorem counts labelled loops on $[n]$ with identity $1$, i.e.\
reduced squares.  Since almost all Latin squares have trivial
autoparatopy group \cite{MW05}, almost all reduced squares have trivial
automorphism group as loops, so the proportion of isomorphism classes
of loops of order $n$ with multiplication group $S_n$ also tends to
$1$ (though with a weaker error bound, that of \cite{MW05}).
\end{remark}

\paragraph{Acknowledgements.}  [Eberhard; Cameron; Wanless.]

\begin{thebibliography}{9}
\bibitem{Cam92} P.~J. Cameron, Almost all quasigroups have rank 2, \emph{Discrete Math.} 106/107 (1992), 111--115.
\bibitem{Cam15} P.~J. Cameron, A niggling problem, blog post, 24 January 2015, with comments by S.~Eberhard, 12--13 October 2016, \url{https://cameroncounts.wordpress.com/2015/01/24/a-niggling-problem/}.
\bibitem{CGW08} N.~J. Cavenagh, C. Greenhill and I.~M. Wanless, The cycle structure of two rows in a random Latin square, \emph{Random Structures Algorithms} 33 (2008), 286--309.
\bibitem{HJ96} R. H\"aggkvist and J.~C.~M. Janssen, All-even Latin squares, \emph{Discrete Math.} 157 (1996), 199--206.
\bibitem{LP93} T. {\L}uczak and L. Pyber, On random generation of the symmetric group, \emph{Combin. Probab. Comput.} 2 (1993), 505--512.
\bibitem{MW05} B.~D. McKay and I.~M. Wanless, On the number of Latin squares, \emph{Ann. Comb.} 9 (2005), 335--344.
\bibitem{PS80} C.~E. Praeger and J. Saxl, On the orders of primitive permutation groups, \emph{Bull. London Math. Soc.} 12 (1980), 303--307.
\bibitem{Pyb93} L. Pyber, Enumerating finite groups of given order, \emph{Ann. of Math.} 137 (1993), 203--220.
\bibitem{RT25} C.~M. Roney-Dougal and G. Tracey, Subgroups of symmetric groups: enumeration and asymptotic properties, arXiv:2503.05416 (2025).
\bibitem{Smith} J.~D.~H. Smith, \emph{An Introduction to Quasigroups and Their Representations}, Chapman \& Hall/CRC, 2007.
\bibitem{vLW} J.~H. van Lint and R.~M. Wilson, \emph{A Course in Combinatorics}, 2nd ed., Cambridge University Press, 2001.
\end{thebibliography}
\end{document}
